\section{Motivation}
\label{section:1.1}

The global online freelance economy has expanded rapidly: in 2023 it was valued at approximately USD\,455 billion, with more than 1.5 billion participants worldwide \cite{statista2024freelance}. Despite their scale, contemporary platforms exhibit four structural deficiencies that impose substantial costs on participants.

\textbf{Intermediary rent.} Every interaction on a centralized platform is mediated by a private corporation that simultaneously custodies escrow, enforces contracts, adjudicates disputes, and sets platform rules. Service fees in the 5--20\% range are routinely levied; Upwork alone collected in excess of USD\,600 million in fee revenue during 2022 \cite{upwork2023}.

\textbf{Payment insecurity.} Platform terms of service typically reserve the right to suspend accounts without notice, freezing held balances. A 2023 survey reports that 34\% of freelancers experienced at least one non-payment incident within twelve months \cite{payoneer2023freelancer}.

\textbf{Reputation lock-in.} Rating systems reside in proprietary databases. A reputation accrued on Upwork cannot be ported to Fiverr or any third-party application, and is susceptible to fake reviews, suppression of negative feedback, and algorithmic bias.

\textbf{Identity and Sybil exposure.} Pseudonymous accounts enable counterfeit identities and review manipulation. Centralized platforms mitigate this through proprietary identity verification, but the resulting biometric data constitutes a centralized repository attractive to adversarial exfiltration.

Public smart-contract blockchains offer an architecturally distinct mechanism for multi-party trust. Smart contracts, first articulated by Szabo \cite{szabo1994smart} and implemented at scale by Ethereum \cite{buterin2014ethereum}, encode contractual obligations as self-executing programs whose behavior does not require a trusted administrator. The Solana blockchain \cite{yakovenko2018solana} further lowers the economic threshold by reporting peak throughput in excess of 50{,}000 \gls{tps} at a typical cost below USD\,0.001 -- compared with Ethereum mainnet fees that have on occasion exceeded USD\,50. These properties motivate the design of a prototype freelance marketplace in which the trust-critical operations are decentralized rather than custodied by a private intermediary.

\section{Objectives and Scope}
\label{section:1.2}

Prior work on decentralized freelancing is fragmentary. Delmolino \emph{et al.}~\cite{delmolino2016step} demonstrated elementary escrow contracts but addressed neither reputation nor identity. Platforms such as Ethlance and Colony provide rudimentary on-chain job posting but, based on their public documentation, do not jointly support biometric verification, milestone escrow, AI-assisted matching, and structured governance within a single system. To the author's knowledge, no extant decentralized system addresses the complete lifecycle of a professional freelance engagement in a unified manner, although a systematic comparison would require dedicated empirical study. The present landscape exhibits four observed gaps: (i)~no widely deployed open-source suite covers the full freelance job lifecycle; (ii)~biometric identity verification with on-chain attestation is uncommon among decentralized platforms; (iii)~SBT reputation, ZK credential proofs, AI risk assessment, and DAO governance are rarely integrated within a single Anchor program; and (iv)~Ethereum-based platforms tend to render small-value engagements economically unattractive at typical gas levels.

This thesis accordingly pursues four objectives, all framed as a proof-of-concept on Solana Devnet:
\begin{enumerate}[noitemsep]
    \item Design and implement an Anchor smart-contract suite covering job management, dual-track escrow, SBT reputation, biometric KYC, DAO governance, social recovery, ZK credential storage, and an AI oracle interface;
    \item Construct a working Next.js front-end supporting both connected and guest browsing;
    \item Integrate browser-side face-recognition inference for a privacy-respecting initial implementation of identity verification;
    \item Demonstrate technical feasibility on Solana Devnet through measured deployment, instruction exercise, and reported costs.
\end{enumerate}

The scope is bounded to a Devnet deployment. Mainnet release, regulatory compliance, third-party security audit, large-scale user studies, full on-chain ZKP verification, and the staking-reward sub-module of the DAO are explicitly identified as out of scope and deferred to future work.

\section{Tentative Solution}
\label{section:1.3}

FreelanceChain adopts a blockchain-first architecture in which the trust-critical operations execute on-chain through Anchor smart contracts. \textbf{Escrow} funds are locked in Program Derived Address (PDA) accounts whose control resides exclusively in contract logic; release requires either client approval, milestone completion, or arbitral ruling. \textbf{Reputation} is minted as non-transferable Soulbound Tokens~\cite{weyl2022decentralized} carrying a numerical rating, comment, and SHA-256 verification hash, which is intended to discourage a secondary market in reputations. \textbf{Biometric KYC} is performed entirely in-browser via \emph{@vladmandic/face-api}: the Euclidean distance between document and selfie face descriptors is compared against a threshold of $0.50$, and only the distance (encoded as basis points) and outcome are written on-chain. \textbf{AI risk assessment} delegates inference to \emph{llama-3.3-70b-versatile} via the Groq API, producing a structured four-axis advisory score and a grounded conversational follow-up; the subsystem is positioned strictly as advisory so that human judgement remains decisive and no AI output is committed on-chain. \textbf{DAO governance} administers platform parameters through reputation-weighted token-stake proposals; the staking-reward sub-module is implemented as a stub and is acknowledged as incomplete.

Solana was selected for its Proof-of-History consensus, which delivers sub-second finality at a fraction of the cost of Ethereum at typical gas levels. Anchor was selected for its strongly typed Rust abstraction, which reduces the smart-contract vulnerability surface relative to hand-written Solana programs.

\section{Thesis Organization}
\label{section:1.4}

Chapter~2 surveys incumbent platforms and derives functional and non-functional requirements. Chapter~3 presents the underlying technologies and the rationale for each selection. Chapter~4 details the system architecture, account schemas, storage design, quantitative measurements, threat model, and testing methodology. Chapter~5 expounds the principal technical contributions: dual-track escrow, SBT reputation, biometric KYC, AI integration, and DAO governance. Chapter~6 concludes with a comparison against existing systems and identifies future research directions.
